%%%%%%%%%%%%%%%%%%%%%%%%%%%%%%%%%%%%%%%%%%%%%%%%%%%%%%%%%%%%%%%%%%%%%%%%%%%%%%%
%% TECHNICAL EVALUATION REPORT -- TEMPLATE
%% ---------------------------------------------------------------------------
%% Reusable skeleton for single-paper technical evaluation reports that are
%% later integrated into a Doctoral Examination Report.
%%
%% Compile with:  pdflatex main.tex  (run twice for the ToC and cross-refs)
%% Requirements:  a standard TeX Live installation; no external .sty files.
%%
%% HOW TO USE
%%   1. Edit the METADATA block below (title, author, version examined, dates).
%%   2. Set \pk to a short key for this paper (used to namespace all labels).
%%   3. Fill the skeleton sections. Commented examples show the intended form.
%%   4. Toggle \verificationtrue / \verificationfalse to include or suppress
%%      the independent-verification appendix, and \draftnotestrue to show
%%      internal notes while drafting.
%%
%% CONVENTIONS ENCODED HERE
%%   [M] = bears on correctness or the substance of a claim.
%%   [R] = bears on presentation, attribution or scope.
%%   Every substantive defect is stated once as a numbered Finding, then
%%   referenced (never restated in full) in the summary, the criticisms list
%%   and the conclusion.
%%%%%%%%%%%%%%%%%%%%%%%%%%%%%%%%%%%%%%%%%%%%%%%%%%%%%%%%%%%%%%%%%%%%%%%%%%%%%%%

\documentclass[11pt]{article}

%%%%%%%%%%%%%%%%%%%%%%%%%%%%%%%%%%% PACKAGES %%%%%%%%%%%%%%%%%%%%%%%%%%%%%%%%%%
\usepackage[utf8]{inputenc}
\usepackage[T1]{fontenc}
\IfFileExists{lmodern.sty}{\usepackage{lmodern}}{} % optional; degrades gracefully
\usepackage[margin=2.6cm]{geometry}
\usepackage{amsmath,amssymb,amsthm}
\usepackage{booktabs}
\usepackage{longtable}
\usepackage{array}
\usepackage{enumitem}
\usepackage{graphicx}
\usepackage[expansion=false]{microtype}
\usepackage{xcolor}
\usepackage{fancyhdr}
\usepackage[colorlinks=true,linkcolor=black,citecolor=black,urlcolor=blue]{hyperref}

%%%%%%%%%%%%%%%%%%%%%%%%%%%%%%%%%% METADATA %%%%%%%%%%%%%%%%%%%%%%%%%%%%%%%%%%%
%% Edit only this block when starting a new report.

\newcommand{\reportKind}{TECHNICAL EVALUATION REPORT}
\newcommand{\paperTitle}{Full Title of the Manuscript Under Review:
                         Subtitle if Any}
\newcommand{\paperShortTitle}{Short Running Title}      % used in the header
\newcommand{\paperAuthor}{A. N. Author}
\newcommand{\paperVersion}{Version examined: \url{http://dx.doi.org/10.xxxx/ssrn.xxxxxxx}
                           (Month Day, Year)}
\newcommand{\evalDate}{Evaluation date: DD Month YYYY}
\newcommand{\evalContext}{Prepared as part of the technical examination of a
  portfolio submitted for the Ph.D. in Mathematical Physics by Published Works}
\newcommand{\reportStatus}{Working technical report for later integration into
  the Doctoral Examination Report}

%% Label namespace for this paper: all findings are labelled \pk:f:<name>,
%% so that reports can be concatenated without label collisions.
\newcommand{\pk}{P0}

%%%%%%%%%%%%%%%%%%%%%%%%%%%%%%%%%%% TOGGLES %%%%%%%%%%%%%%%%%%%%%%%%%%%%%%%%%%%

\newif\ifverification   \verificationtrue   % include the verification appendix
\newif\ifclassification \classificationtrue % include the result-by-result table
\newif\ifdraftnotes     \draftnotesfalse    % show internal drafting notes

%%%%%%%%%%%%%%%%%%%%%%%%%%%%%%%%%%% MACROS %%%%%%%%%%%%%%%%%%%%%%%%%%%%%%%%%%%%

%% Severity markers.
\newcommand{\defM}{\textbf{[M]}}   % material: correctness or substance
\newcommand{\defR}{\textbf{[R]}}   % remark: presentation, attribution, scope

%% Findings: state once, reference everywhere else.
\theoremstyle{definition}
\newtheorem{finding}{Finding}
\newtheorem*{remarkE}{Remark}

\newcommand{\flabel}[1]{\label{\pk:f:#1}}          % inside a finding
\newcommand{\fref}[1]{Finding~\ref{\pk:f:#1}}      % anywhere else
\newcommand{\fpref}[1]{(\fref{#1})}

%% Internal notes, invisible unless \draftnotestrue.
\newcommand{\note}[1]{\ifdraftnotes\textcolor{red}{\small[\,#1\,]}\fi}

%% Shorthands commonly needed in this series. Extend per paper.
\newcommand{\op}{\mathrm{op}}
\newcommand{\RR}{\mathbb{R}}
\newcommand{\PP}{\mathbb{P}}
\newcommand{\EE}{\mathbb{E}}
\newcommand{\norm}[1]{\left\lVert #1 \right\rVert}
\newcommand{\abs}[1]{\left\lvert #1 \right\rvert}
\newcommand{\inner}[2]{\langle #1 , #2 \rangle}

%% "\S" references always point at the manuscript under review, never at this
%% report; this macro makes that explicit and greppable.
\newcommand{\ms}[1]{\S#1}

%%%%%%%%%%%%%%%%%%%%%%%%%%%%%%%%%% PAGE STYLE %%%%%%%%%%%%%%%%%%%%%%%%%%%%%%%%%

\pagestyle{fancy}
\fancyhf{}
\lhead{\small Technical Evaluation}
\fancyhead[R]{\small\itshape \paperShortTitle}
\fancyfoot[C]{\thepage}
\renewcommand{\headrulewidth}{0.3pt}

\begin{document}

%%%%%%%%%%%%%%%%%%%%%%%%%%%%%%%%% TITLE PAGE %%%%%%%%%%%%%%%%%%%%%%%%%%%%%%%%%%

\begin{titlepage}
\centering
\vspace*{2cm}
{\LARGE\bfseries \reportKind\par}
\vspace{1.2cm}
{\Large\bfseries \paperTitle\par}
\vspace{0.7cm}
{\large \paperAuthor\par}
\vspace{1.4cm}
{\normalsize \paperVersion\par}
\vspace{1cm}
{\itshape \evalContext\par}
\vfill
{\evalDate\par}
\vspace{0.8cm}
{\bfseries Status: \reportStatus\par}
\end{titlepage}

\tableofcontents
\newpage

%%%%%%%%%%%%%%%%%%%%%%%%%%%%%%%%%%%%%%%%%%%%%%%%%%%%%%%%%%%%%%%%%%%%%%%%%%%%%%%
\section{Summary of the submission}
%%%%%%%%%%%%%%%%%%%%%%%%%%%%%%%%%%%%%%%%%%%%%%%%%%%%%%%%%%%%%%%%%%%%%%%%%%%%%%%
%% One paragraph: the question the paper asks and the form of its answer.
%% Then one paragraph tracing the architecture section by section, so that a
%% reader who has not read the manuscript can follow the rest of this report.
%% Close by naming the evaluation standard actually applied and what was done
%% to test the paper against it.

The paper asks \note{question} and answers with \note{form of the answer}.

The argument proceeds in \note{n} stages. \ms{2} fixes the setting: \note{\dots}
\ms{3} \note{\dots} \ms{4} \note{\dots} \ms{5} \note{\dots}

The primary evaluation criteria applied here are logical consistency with a
testable hypothesis and empirical reproducibility. \note{State what the paper
claims about its own verifiability, and what was done to test that claim ---
independent re-derivation, exact arithmetic, simulation replication, data
re-analysis, as applicable.}

%%%%%%%%%%%%%%%%%%%%%%%%%%%%%%%%%%%%%%%%%%%%%%%%%%%%%%%%%%%%%%%%%%%%%%%%%%%%%%%
\section{Overall assessment}
%%%%%%%%%%%%%%%%%%%%%%%%%%%%%%%%%%%%%%%%%%%%%%%%%%%%%%%%%%%%%%%%%%%%%%%%%%%%%%%
%% Structure enforced here, in this order:
%%   (a) verdict on correctness and reproducibility, with the evidence;
%%   (b) what the paper does well, itemised, before any criticism;
%%   (c) the material defects, one short paragraph each, each pointing at the
%%       Finding that states it in full;
%%   (d) the minor items, as a list.
%% Nothing in this section should be argued for the first time: it is a map.

\paragraph{Correctness and reproducibility.}
\note{Verdict in one sentence, then the evidence.} The following components
were recomputed and confirmed:
\begin{itemize}
  \item \note{quantity, with its exact value}
  \item \note{quantity, with its exact value}
\end{itemize}
\note{Say explicitly what kind of check was performed --- exact rational
arithmetic, floating point with a stated tolerance, symbolic re-derivation, or
by hand --- and what a zero residual does and does not establish.}

\paragraph{Features that deserve notice.}
\begin{itemize}
  \item \textit{\note{Label}.} \note{One or two sentences.}
  \item \textit{\note{Label}.} \note{One or two sentences.}
\end{itemize}

\paragraph{Material defects.}
The following discrepancies between what is proved and what is announced are
recorded. \note{State whether any touches a theorem.}

\note{One short paragraph per material finding, each ending with a pointer:}
\fpref{key-one}.

\paragraph{Minor items.}
\begin{itemize}
  \item \note{item} \fpref{key-two}
  \item \note{item}
\end{itemize}

%%%%%%%%%%%%%%%%%%%%%%%%%%%%%%%%%%%%%%%%%%%%%%%%%%%%%%%%%%%%%%%%%%%%%%%%%%%%%%%
\section{Section-by-section analysis}
%%%%%%%%%%%%%%%%%%%%%%%%%%%%%%%%%%%%%%%%%%%%%%%%%%%%%%%%%%%%%%%%%%%%%%%%%%%%%%%
%% One subsection per section of the manuscript. Within each, use
%% \paragraph{...} headings that state the verdict on a single result, so the
%% reader can scan the verdicts without reading the arguments. Findings are
%% raised here, at the point where the evidence for them appears.

\subsection{Setting and definitions (\ms{2})}

\paragraph{Proposition 2.2 is correct.} \note{Reproduce the step that matters,
not the whole proof. Say which identities the proof uses and why they hold.}

\paragraph{Lemma 2.3 is correct.} \note{Verdict, then the verification: what was
recomputed, on what instances, with what residual.}

%% --- Example of a fully formed finding; delete or adapt. ---------------------
\begin{finding}[One-line statement of the defect, not of the topic]\flabel{key-one}
\note{Quote or paraphrase precisely what the manuscript claims, with its
location. Then state what is actually the case, with the evidence. Then state
the consequence: what breaks, what does not, and what the repair is. A finding
that does not end with a repair is incomplete.}
\defM
\end{finding}
%% ----------------------------------------------------------------------------

\subsection{Main construction (\ms{3}--\ms{4})}

\paragraph{Theorem 4.2 is correct, and is an identity.} \note{\dots}

\subsection{Interpretation and connections (\ms{5})}

\paragraph{\note{Claim made in passing}.} \note{Claims made in passing are held
to the same standard as theorems, and are the usual source of [R] items.}
\defR

\subsection{Sharpness and limiting cases (\ms{6})}

\begin{finding}[\note{Statement}]\flabel{key-two}
\note{\dots}
\defR
\end{finding}

%%%%%%%%%%%%%%%%%%%%%%%%%%%%%%%%%%%%%%%%%%%%%%%%%%%%%%%%%%%%%%%%%%%%%%%%%%%%%%%
\section{Originality: what is new and what is standard}
%%%%%%%%%%%%%%%%%%%%%%%%%%%%%%%%%%%%%%%%%%%%%%%%%%%%%%%%%%%%%%%%%%%%%%%%%%%%%%%
%% Three paragraphs, always in this order. The first protects the author from
%% the charge of claiming known material; the second states the contribution in
%% the evaluator's own words; the third places the work on a level.

\paragraph{Standard, and correctly attributed.} \note{Enumerate the inherited
apparatus and say whether each is cited at its point of use.}

\paragraph{The paper's own.} \note{Enumerate the contributions that survive the
audit, one clause each, and say for each what makes it non-obvious.}

\paragraph{Level.} \note{What kind of work this is, what machinery it requires,
and where its originality actually lies.}

%%%%%%%%%%%%%%%%%%%%%%%%%%%%%%%%%%%%%%%%%%%%%%%%%%%%%%%%%%%%%%%%%%%%%%%%%%%%%%%
\section{Priority in the literature}
%%%%%%%%%%%%%%%%%%%%%%%%%%%%%%%%%%%%%%%%%%%%%%%%%%%%%%%%%%%%%%%%%%%%%%%%%%%%%%%
%% Antecedents the manuscript does not cite, each with a full reference, a
%% statement of what is genuinely shared, and a statement of what remains the
%% paper's own. Positioning items sharpen a claim; they should never be written
%% as accusations of overlap.

\begin{finding}[\note{n} positioning items are missing]\flabel{priority}
\emph{\note{Antecedent one}.} \note{Author, title, journal, volume, year, pages.
What is shared. What is different. Why making the comparison would sharpen
rather than dilute the claim.}

\emph{\note{Antecedent two}.} \note{\dots}
\defR
\end{finding}

% %%%%%%%%%%%%%%%%%%%%%%%%%%%%%%%%%%%%%%%%%%%%%%%%%%%%%%%%%%%%%%%%%%%%%%%%%%%%%%%
% \ifclassification
% \section{Result-by-result classification}
% %%%%%%%%%%%%%%%%%%%%%%%%%%%%%%%%%%%%%%%%%%%%%%%%%%%%%%%%%%%%%%%%%%%%%%%%%%%%%%%
% %% One row per component of the manuscript. Weight is the component's load in
% %% the paper's own architecture (Major / Substantial / Moderate / Minor /
% %% Supporting), not the severity of any defect.

% {\small
% \begin{longtable}{@{}p{3.0cm}p{2.5cm}p{2.6cm}p{1.9cm}p{3.0cm}@{}}
% \toprule
% \textbf{Component} & \textbf{Correctness} & \textbf{Originality} &
% \textbf{Weight} & \textbf{Recommended action}\\
% \midrule
% \endfirsthead
% \toprule
% \textbf{Component} & \textbf{Correctness} & \textbf{Originality} &
% \textbf{Weight} & \textbf{Recommended action}\\
% \midrule
% \endhead
% \bottomrule
% \endfoot
% \bottomrule
% \endlastfoot
% \note{Definition / metric} & \textbf{Correct; load-bearing} & The paper's own &
% \textbf{Major} & \note{Action} \fpref{key-one}\\[0.3em]
% \note{Lemma 2.3} & \textbf{Correct; verified} & Moderate & Substantial & None
% material\\[0.3em]
% \note{Claim made in passing} & Unsupported as stated & --- & Minor & Present as
% an analogy\\[0.3em]
% Reproducibility & \note{verdict} & --- & \textbf{Major} & \note{action}\\
% \end{longtable}}
% \fi

%%%%%%%%%%%%%%%%%%%%%%%%%%%%%%%%%%%%%%%%%%%%%%%%%%%%%%%%%%%%%%%%%%%%%%%%%%%%%%%
\section{Specific criticisms}
%%%%%%%%%%%%%%%%%%%%%%%%%%%%%%%%%%%%%%%%%%%%%%%%%%%%%%%%%%%%%%%%%%%%%%%%%%%%%%%
%% The actionable list, [M] items first. Each entry is one sentence of defect
%% plus one clause of repair, with a pointer to the Finding. Nothing new is
%% argued here.

Items are marked \defM\ where they bear on the correctness or the substance of a
claim, and \defR\ where they concern presentation, attribution or scope.

\begin{description}[leftmargin=2.4em,style=nextline]
\item[\defM\ \note{One-line statement of the defect}]
  \fpref{key-one}. \note{Repair in one clause.}
\item[\defR\ \note{One-line statement of the defect}]
  \fpref{key-two}. \note{Repair in one clause.}
\item[\defR\ \note{Positioning}] \fpref{priority}.
\end{description}

%%%%%%%%%%%%%%%%%%%%%%%%%%%%%%%%%%%%%%%%%%%%%%%%%%%%%%%%%%%%%%%%%%%%%%%%%%%%%%%
\section{Questions to the author}
%%%%%%%%%%%%%%%%%%%%%%%%%%%%%%%%%%%%%%%%%%%%%%%%%%%%%%%%%%%%%%%%%%%%%%%%%%%%%%%
%% Written for a viva: each question is answerable, points at a specific object
%% in the manuscript, and carries the evaluator's own computation where one was
%% made. Order: material defects first, then scope, then open problems.

\begin{enumerate}[leftmargin=2.2em]
\item \note{Question bearing on a material finding, with the evaluator's
      counter-computation stated as a number.}
\item \note{Question on scope or on a hypothesis silently dropped.}
\item \note{Question inviting the author to extend the work at a point the
      audit identified as tractable.}
\end{enumerate}

%%%%%%%%%%%%%%%%%%%%%%%%%%%%%%%%%%%%%%%%%%%%%%%%%%%%%%%%%%%%%%%%%%%%%%%%%%%%%%%
\section{Technical conclusion}
%%%%%%%%%%%%%%%%%%%%%%%%%%%%%%%%%%%%%%%%%%%%%%%%%%%%%%%%%%%%%%%%%%%%%%%%%%%%%%%
%% Four paragraphs: what was verified and how; what the architecture achieves;
%% the material defects gathered in one place with their weight stated; the
%% residual matters of scope. The conclusion must be readable standing alone,
%% because it is what gets carried into the integrative report.

\note{Paragraph 1 --- verification and reproducibility.}

\note{Paragraph 2 --- the architecture and what is genuinely new.}

\note{Paragraph 3 --- the material defects, and whether any touches a theorem.}

\note{Paragraph 4 --- matters of scope, positioning and terminology.}

% %%%%%%%%%%%%%%%%%%%%%%%%%%%%%%%%%%%%%%%%%%%%%%%%%%%%%%%%%%%%%%%%%%%%%%%%%%%%%%%
% \ifverification
% \appendix
% \section{Independent verification record}
% %%%%%%%%%%%%%%%%%%%%%%%%%%%%%%%%%%%%%%%%%%%%%%%%%%%%%%%%%%%%%%%%%%%%%%%%%%%%%%%
% %% The audit trail. State the arithmetic regime once, then one row per claim
% %% checked, including the checks that returned a negative result --- those are
% %% the rows that carry the findings.

% All computations were performed from the printed definitions alone.

% \begin{center}
% {\small
% \begin{longtable}{@{}p{6.0cm}p{4.0cm}p{4.3cm}@{}}
% \toprule
% \textbf{Claim} & \textbf{Method} & \textbf{Outcome}\\
% \midrule
% \endfirsthead
% \toprule
% \textbf{Claim} & \textbf{Method} & \textbf{Outcome}\\
% \midrule
% \endhead
% \bottomrule
% \endfoot
% \bottomrule
% \endlastfoot
% \note{Example 3.2: displayed matrices} & Exact & As printed\\[0.35em]
% \note{Lemma 2.3: isometry} & \note{12 random rational families} & Residual
% exactly $0$\\[0.35em]
% \note{Counter-check that produced a finding} & \note{method} &
% \textbf{No}: \note{discrepancy exactly \dots}\\
% \end{longtable}}
% \end{center}

% \subsection{\note{Quantified sub-check}}
% %% Optional: a short table isolating one quantitative claim of the audit, e.g.
% %% how badly an identity fails under an alternative hypothesis.

% \begin{center}
% {\small
% \begin{tabular}{@{}lccc@{}}
% \toprule
% \textbf{Trial} & \textbf{\note{regime A}} & \textbf{\note{regime B}} &
% \textbf{relative gap}\\
% \midrule
% 1 & \note{value} & \note{value} & \note{\%}\\
% 2 & \note{value} & \note{value} & \note{\%}\\
% \bottomrule
% \end{tabular}}
% \end{center}
% \fi

\end{document}
